% ===== PRÉAMBULE CANONIQUE — à coller VERBATIM en tête de CHAQUE .tex =====
\documentclass[11pt,a4paper]{article}
\usepackage[utf8]{inputenc}\usepackage[T1]{fontenc}\usepackage{lmodern}
\usepackage[french]{babel}
\usepackage{amsmath,amssymb,mathtools}\usepackage{icomma}
\usepackage{siunitx}\sisetup{locale=FR}  % \qty{}{}, \si{}, \ang{}
\usepackage{xcolor}\usepackage{colortbl}
\usepackage{tikz}\usepackage{pgfplots}\pgfplotsset{compat=1.18}
\usepackage[siunitx,europeanresistors]{circuitikz} % circuits (élec.)
\usepackage[most]{tcolorbox}\usepackage{enumitem}\usepackage{array,booktabs}
\usepackage[a4paper,margin=2.2cm]{geometry}
\usetikzlibrary{arrows.meta,calc,angles,quotes,positioning,patterns,%
  backgrounds,shapes.geometric,decorations.pathmorphing,decorations.markings,intersections}
\definecolor{primary}{RGB}{222,49,99}\definecolor{primarydark}{RGB}{150,22,55}
\definecolor{defcol}{RGB}{0,114,178}\definecolor{methcol}{RGB}{52,92,148}
\definecolor{excol}{RGB}{0,135,90}\definecolor{attncol}{RGB}{200,40,55}
\tcbset{boxbase/.style={enhanced,breakable,boxrule=0.9pt,arc=2.5pt,
  left=3mm,right=3mm,top=2.3mm,bottom=2.3mm,fonttitle=\bfseries,coltitle=white}}
% --- boîtes TUTORIEL (noms EXACTS, ne pas renommer) ---
\newtcolorbox{cle}[1][]{boxbase,colback=defcol!6,colframe=defcol,
  title={\textsc{À retenir}\ifx\relax#1\relax\relax\else\textmd{ : \itshape #1}\fi}}
\newtcolorbox{methode}[1][]{boxbase,colback=methcol!7,colframe=methcol,sharp corners,
  title={\textsc{Méthode}\ifx\relax#1\relax\relax\else\textmd{ --- \itshape #1}\fi}}
\newtcolorbox{exemplebox}[1][]{boxbase,colback=excol!5,colframe=excol,
  title={\textsc{Exemple}\ifx\relax#1\relax\relax\else\textmd{ : \itshape #1}\fi}}
\newtcolorbox{piege}[1][]{boxbase,colback=attncol!6,colframe=attncol,
  title={\textsc{Piège}\ifx\relax#1\relax\relax\else\textmd{ : \itshape #1}\fi}}
\newtcolorbox{remarque}[1][]{boxbase,colback=black!4,colframe=black!45,
  title={\textsc{Remarque}\ifx\relax#1\relax\relax\else\textmd{ : \itshape #1}\fi}}
% --- boîtes EXERCICE / CORRIGÉ (noms EXACTS, auto-comptées) ---
\newtcolorbox[auto counter]{exo}[1][]{boxbase,colback=primary!4,colframe=primary,
  title={\textsc{Exercice}~\thetcbcounter\ifx\relax#1\relax\relax\else\textmd{ --- \itshape #1}\fi}}
\newtcolorbox[auto counter]{corrige}[1][]{boxbase,colback=excol!4,colframe=excol,
  title={\textsc{Corrigé}~\thetcbcounter\ifx\relax#1\relax\relax\else\textmd{ --- \itshape #1}\fi}}
\newcommand{\vv}[1]{\vec{#1}}\newcommand{\ds}{\displaystyle}
\newcommand{\R}{\mathbb{R}}\newcommand{\N}{\mathbb{N}}\newcommand{\Z}{\mathbb{Z}}
% =========================================================================
\begin{document}
% THEME_TITLE: Plan incliné, réaction normale et frottement
\sisetup{per-mode=symbol}

\noindent Sur une surface, un corps subit deux forces de contact : la \textbf{réaction normale} $\vv N$
(perpendiculaire) et le \textbf{frottement} $\vv F_f$ (parallèle). Le \textbf{plan incliné} est
l'application reine des trois lois.

\begin{cle}[plan incliné : les résultats à connaître]
On choisit un repère \emph{aligné sur la pente} ($Ox$ le long du plan, $Oy$ perpendiculaire).
On projette le poids, puis on écrit le PFD axe par axe :
\[ G_x=mg\sin\alpha \quad(\text{le long du plan}),\qquad G_y=mg\cos\alpha \quad(\perp\ \text{au plan}), \]
\[ \text{Selon } Oy:\ \boxed{N=mg\cos\alpha}\qquad\qquad \text{Selon } Ox:\ \boxed{G_x-F_f=ma}. \]
Sur un \emph{sol horizontal} ($\alpha=0$) : $N=mg$. \quad Sans frottement : $a=g\sin\alpha$.
\end{cle}

\begin{center}
\begin{tikzpicture}[>=Stealth,scale=1.05]
  \coordinate (O) at (0,0); \coordinate (B) at (6,0); \coordinate (C) at (6,3.464);
  \draw[thick,fill=black!5] (O)--(B)--(C)--cycle;
  \draw (0.95,0) arc (0:30:0.95); \node at (1.3,0.24){$\alpha$};
  \coordinate (P) at (3,1.732);
  \fill[defcol] (P) circle (0.13);
  \draw[->,line width=1.1pt,attncol] (P) -- ++(0,-1.8) node[below]{$\vv G=m\vv g$};
  \draw[->,line width=1pt,attncol!65] (P) -- ++(-0.78,-0.45) node[left=0pt]{$\vv G_x$};
  \draw[->,line width=1pt,attncol!65] (P) -- ++(0.78,-1.35) node[right=0pt]{$\vv G_y$};
  \draw[dashed,black!40] ($(P)+(-0.78,-0.45)$) -- ($(P)+(0,-1.8)$);
  \draw[dashed,black!40] ($(P)+(0.78,-1.35)$) -- ($(P)+(0,-1.8)$);
  \draw[->,line width=1.1pt,excol] (P) -- ++(-0.7,1.212) node[above left]{$\vv N$};
  \draw[->,line width=1.1pt,methcol] (P) -- ++(0.952,0.55) node[above right]{$\vv F_f$};
  \draw[->,black!55] (1.15,1.35) -- ++(30:0.75) node[right,font=\scriptsize]{$x$};
  \draw[->,black!55] (1.15,1.35) -- ++(120:0.75) node[left,font=\scriptsize]{$y$};
  \node[anchor=west,align=left,font=\small] at (6.35,2.55)
     {$G_x=mg\sin\alpha$\\[3pt]$G_y=mg\cos\alpha$};
  \node[anchor=west,align=left,font=\small] at (6.35,0.95)
     {\textcolor{primarydark}{$Oy$ :}\ $N=G_y$\\[3pt]\textcolor{primarydark}{$Ox$ :}\ $G_x-F_f=ma$};
\end{tikzpicture}
\end{center}

\begin{methode}[le frottement selon l'état du corps]
\begin{itemize}[leftmargin=1.5em,topsep=2pt,itemsep=1.5pt]
  \item \textbf{Corps immobile} (frottement \emph{statique}) : $F_f$ \emph{s'ajuste} pour équilibrer la
        force appliquée, jusqu'à un maximum : $F_f\le \mu_s N$.
  \item \textbf{Corps qui glisse} (frottement \emph{cinétique}) : $F_f=\mu_c N$ (valeur fixée).
  \item \textbf{À vitesse constante} (MRU) : $a=0$, donc le frottement \emph{égale} la force motrice
        (équilibre) — inutile de connaître $\mu$ !
\end{itemize}
\end{methode}

\begin{exemplebox}[cycliste sur une pente de \ang{30}]
$G=\qty{802}{\newton}$, $F_f=\qty{20}{\newton}$, $g=\qty{10}{\metre\per\second\squared}$, donc
$m=\tfrac{G}{g}=\qty{80,2}{\kilogram}$. Composante motrice $G_x=G\sin\ang{30}=\qty{401}{\newton}$.
Selon $Ox$ : $a=\dfrac{G_x-F_f}{m}=\dfrac{401-20}{\num{80,2}}\approx\qty{4,75}{\metre\per\second\squared}$.
\end{exemplebox}

\begin{piege}[les pièges classiques]
\begin{itemize}[leftmargin=1.4em,topsep=2pt,itemsep=1.5pt]
  \item À \textbf{vitesse constante}, $F_f=$ force appliquée ; $\mu_c N$ n'est que la valeur \emph{maximale}
        possible — le coefficient $\mu$ donné est souvent un \emph{piège} (donnée inutile).
  \item Sur un plan incliné, $N=mg\cos\alpha$ (et \emph{non} $mg$) : la normale ne compense que $G_y$.
  \item Sans frottement, $a=g\sin\alpha$ \textbf{ne dépend pas de la masse} : deux corps de masses
        différentes descendent identiquement.
  \item Composante \emph{motrice} $=mg\sin\alpha$ (avec le \textbf{sinus}) ; composante plaquante $=mg\cos\alpha$.
\end{itemize}
\end{piege}
\end{document}
